\documentclass[10pt,a4paper]{amsart}
\usepackage{amsmath,amssymb,amsthm,graphicx}
\usepackage{verbatim}
\newenvironment{reference}{\comment}{\endcomment}
%\newenvironment{reference}{}{}
\newenvironment{slogan}{\comment}{\endcomment}
\newenvironment{history}{\comment}{\endcomment}

% For commutative diagrams we use Xy-pic
\usepackage[all]{xy}

% We use 2cell for 2-commutative diagrams.
\xyoption{2cell}
% Theorem environments.
%
\theoremstyle{plain}
\newtheorem{theorem}[subsection]{Theorem}
\newtheorem{proposition}[subsection]{Proposition}
\newtheorem{lemma}[subsection]{Lemma}

\theoremstyle{definition}
\newtheorem{definition}[subsection]{Definition}
\newtheorem{example}[subsection]{Example}
\newtheorem{exercise}[subsection]{Exercise}
\newtheorem{situation}[subsection]{Situation}

\theoremstyle{remark}
\newtheorem{remark}[subsection]{Remark}
\newtheorem{remarks}[subsection]{Remarks}

\numberwithin{equation}{subsection}

% Macros
%
\def\lim{\mathop{\mathrm{lim}}\nolimits}
\def\colim{\mathop{\mathrm{colim}}\nolimits}
\def\Spec{\mathop{\mathrm{Spec}}}
\def\Hom{\mathop{\mathrm{Hom}}\nolimits}
\def\Ext{\mathop{\mathrm{Ext}}\nolimits}
\def\SheafHom{\mathop{\mathcal{H}\!\mathit{om}}\nolimits}
\def\SheafExt{\mathop{\mathcal{E}\!\mathit{xt}}\nolimits}
\def\Sch{\mathit{Sch}}
\def\Mor{\mathop{\mathrm{Mor}}\nolimits}
\def\Ob{\mathop{\mathrm{Ob}}\nolimits}
\def\Sh{\mathop{\mathit{Sh}}\nolimits}
\def\NL{\mathop{N\!L}\nolimits}
\def\CH{\mathop{\mathrm{CH}}\nolimits}
\def\proetale{{pro\text{-}\acute{e}tale}}
\def\etale{{\acute{e}tale}}
\def\QCoh{\mathit{QCoh}}
\def\Ker{\mathop{\mathrm{Ker}}}
\def\Im{\mathop{\mathrm{Im}}}
\def\Coker{\mathop{\mathrm{Coker}}}
\def\Coim{\mathop{\mathrm{Coim}}}

% Boxtimes
%
\DeclareMathSymbol{\boxtimes}{\mathbin}{AMSa}{"02}

%
% Macros for moduli stacks/spaces
%
\def\QCohstack{\mathcal{QC}\!\mathit{oh}}
\def\Cohstack{\mathcal{C}\!\mathit{oh}}
\def\Spacesstack{\mathcal{S}\!\mathit{paces}}
\def\Quotfunctor{\mathrm{Quot}}
\def\Hilbfunctor{\mathrm{Hilb}}
\def\Curvesstack{\mathcal{C}\!\mathit{urves}}
\def\Polarizedstack{\mathcal{P}\!\mathit{olarized}}
\def\Complexesstack{\mathcal{C}\!\mathit{omplexes}}
% \Pic is the operator that assigns to X its picard group, usage \Pic(X)
% \Picardstack_{X/B} denotes the Picard stack of X over B
% \Picardfunctor_{X/B} denotes the Picard functor of X over B
\def\Pic{\mathop{\mathrm{Pic}}\nolimits}
\def\Picardstack{\mathcal{P}\!\mathit{ic}}
\def\Picardfunctor{\mathrm{Pic}}
\def\Deformationcategory{\mathcal{D}\!\mathit{ef}}
\usepackage{cleveref}
\crefname{originallemma}{original lemma}{original lemmas}
\Crefname{originallemma}{Original lemma}{Original lemmas}
\makeatletter
\newlabel{algebra-lemma-directed-colimit-exact}{{00DB}{source}{}{sourcepointer-1}{}}
\newlabel{algebra-lemma-directed-colimit-exact@cref}{{[originallemma][0][]00DB}{[1][0][]source}{}{}{}}
\makeatother
\crefname{originalreference}{original reference}{original references}
\Crefname{originalreference}{Original reference}{Original references}
\makeatletter
\newlabel{lemma-derived-bc}{{09LS}{source}{}{sourcepointer-2}{}}
\newlabel{lemma-derived-bc@cref}{{[originalreference][0][]09LS}{[1][0][]source}{}{}{}}
\makeatother
\title{Derived tensor functors are functorial in the bimodule, and quasi-isomorphisms give isomorphisms of functors}
\author{The Stacks Project Authors}
\date{stacks-project:dga.tex:local-snapshot-20261012; GNU FDL 1.2 or later}
\begin{document}
\maketitle
\noindent\textit{Editorial note.} All delivered mathematical text is the author's own, copied from the pinned Stacks Project chapter file sources/stacks-dga/pinned-original-source.tex; no statement or proof step was rewritten, reordered or removed.\par\medskip
\noindent\textit{Editorial note.} The counted proof is the author's complete argument: the property-(P) assembly of the transformation, the filtered-colimit reduction, the induction on the filtration with the split-inclusion short exact sequences, and the final graded-piece evaluation are all retained verbatim.\par\medskip
\noindent\textit{Editorial note.} The label lemma-functoriality-bc names the counted statement, and the proof environment immediately following it in the source is delivered verbatim as its proof.\par\medskip
\noindent\textit{Editorial note.} The counted statement is carried by the author's own lemma environment in the chapter source.\par\medskip
\noindent\textit{Allowed standard prerequisites (editorial).}\begin{itemize}
\item Computation of derived tensor by property (P) modules. Same chapter, Lemma lemma-derived-bc. Identifies the derived tensors - tensor\_A\textasciicircum{}L N and - tensor\_A\textasciicircum{}L N' with plain tensors on property-(P) modules, the starting point for the transformation and its quasi-isomorphism analysis.
\item Exactness of filtered colimits and the filtration of property-(P) modules. Algebra, Lemma algebra-lemma-directed-colimit-exact; Same chapter, construction of property (P). Lets the tensor with M be read through the filtered system of its filtration stages, whose split inclusions give the short exact sequences the induction runs on.
\end{itemize}\medskip
\ifdefined\hypertarget\hypertarget{sourcepointer-1}{}\fi
\noindent\textit{Source reference.} Original lemma 00DB: sources/stacks-chow/tags.txt, source lines 491--491. Prerequisite pointer invoked at proof line(s) 5592; it names a standard fact of a sibling chapter that the counted argument uses, and every step of the counted proof built on it is retained verbatim in the delivered body.\par\medskip
\ifdefined\hypertarget\hypertarget{sourcepointer-2}{}\fi
\noindent\textit{Source reference.} Original reference 09LS: sources/stacks-dga/pinned-original-source.tex, source lines 5541--5541. Prerequisite pointer invoked at proof line(s) 5581; it names a standard fact of the same chapter that the counted argument uses, and every step of the counted proof built on it is retained verbatim in the delivered body.\par\medskip
\begin{lemma}
\label{lemma-functoriality-bc}
Let $R$ be a ring. Let $(A, \text{d})$ and $(B, \text{d})$ be
differential graded $R$-algebras. Let $f : N \to N'$ be a
homomorphism of differential graded $(A, B)$-bimodules.
Then $f$ induces a morphism of functors
$$
1\otimes f :
- \otimes_A^\mathbf{L} N
\longrightarrow
- \otimes_A^\mathbf{L} N'
$$
If $f$ is a quasi-isomorphism, then $1 \otimes f$ is an isomorphism of
functors.
\end{lemma}

\begin{proof}
Let $M$ be a differential graded $A$-module with property (P).
Then $1 \otimes f : M \otimes_A N \to M \otimes_A N'$ is a
map of differential graded $B$-modules. Moreover, this is functorial
with respect to $M$. Since the functors
$- \otimes_A^\mathbf{L} N$ and $- \otimes_A^\mathbf{L} N'$ are
computed by tensoring on objects with property (P)
(Lemma \ref{lemma-derived-bc}) we obtain a transformation of functors
as indicated.

\medskip\noindent
Assume that $f$ is a quasi-isomorphism. Let $F_\bullet$ be the
given filtration on $M$. Observe that
$M \otimes_A N = \colim F_i(M) \otimes_A N$ and
$M \otimes_A N' = \colim F_i(M) \otimes_A N'$.
Hence it suffices to show that
$F_n(M) \otimes_A N \to F_n(M) \otimes_A N'$
is a quasi-isomorphism (filtered colimits are exact, see
Algebra, Lemma \ref{algebra-lemma-directed-colimit-exact}).
Since the inclusions $F_n(M) \to F_{n + 1}(M)$
are split as maps of graded $A$-modules we see that
$$
0 \to F_n(M) \otimes_A N \to F_{n + 1}(M) \otimes_A N \to
F_{n + 1}(M)/F_n(M) \otimes_A N \to 0
$$
is a short exact sequence of differential graded $B$-modules.
There is a similar sequence for $N'$ and $f$ induces a map
of short exact sequences. Hence by induction on $n$ (starting with $n = -1$
when $F_{-1}(M) = 0$) we conclude that it suffices to show that
the map $F_{n + 1}(M)/F_n(M) \otimes_A N \to F_{n + 1}(M)/F_n(M) \otimes_A N'$
is a quasi-isomorphism. This is true because $F_{n + 1}(M)/F_n(M)$
is a direct sum of shifts of $A$ and the result is true for $A[k]$
as $f : N \to N'$ is a quasi-isomorphism.
\end{proof}

\end{document}
